% !PRINTED 95
\weberpage{95}
%<<CONT
zu setzen, und demnach wird [nach (6), (7), (10)]
\[
\begin{gathered}
\omega'_1 = \frac{a+b\omega}{n}, \quad \omega'_2 = \frac{c+\partial\omega}{n}, \\
a'_1 = m'b, \quad a'_2 = m'\partial.
\end{gathered}
\]
Die Funktion $T'(u, \omega'_1, \omega'_2)$ genügt also den Bedingungen (1):
\[
\begin{gathered}
T(u+\omega'_1) = (-1)^{g'_1} e^{-\pi i m'b(2u+\omega'_1)} T(u), \\
T(u+\omega'_2) = (-1)^{g'_2} e^{-\pi i m'\partial(2u+\omega'_2)} T(u),
\end{gathered}
\]
% [sic] The source prints unprimed T on both rows although the defined function is T'.
und daraus ergibt sich, daß das Produkt
\[
e^{\frac{\pi i m'bu^2}{\omega'_1}} T'(u, \omega'_1, \omega'_2)
\]
eine $\Theta$-Funktion der Ordnung $m'$ ist, mit den Argumenten
\[
\frac{u}{\omega'_1}, \quad \frac{\omega'_2}{\omega'_1}
\]
und der Charakteristik $(g'_1, g'_2)$. Wir können dies in der
Gleichung ausdrücken:
\weq{12}{e^{-\frac{\pi i m' n b u^2}{a+b\omega}} \Theta^{(m')}_{g'_1,g'_2}\left(\frac{nu}{a+b\omega},\ \frac{c+\partial\omega}{a+b\omega}\right) = \Theta^{(m'n)}_{g_1,g_2}(u,\omega),}
worin die Charakteristiken durch (11) bestimmt sind. Die Mittel
zur Darstellung dieser Funktionen sind in § 21 enthalten.

Wir bezeichnen die Transformation von $T$ und $T'$ durch
einen einzelnen Buchstaben $S$ oder durch $(T', T)$, also:
\[
S = (T', T).
\]

Bedeutet $S'$ eine zweite Transformation, durch die $T'$ in $T''$
übergeht, also
\[
S' = (T'', T'),
\]
so können wir daraus eine neue Transformation $S''$ ableiten,
durch die $T$ in $T''$ übergeht. Diese heißt aus $S$ und $S'$ zusammengesetzt
und wird so bezeichnet:
\[
S'' = S'S
\]
oder
\[
(T'', T) = (T'', T')(T', T).
\]
Bei dieser Zusammensetzung gilt im allgemeinen nicht das
kommutative Gesetz; es ist also $SS'$ von $S'S$ verschieden. Es
gilt aber das assoziative Gesetz, das sich in der Formel ausspricht:
\[
(T''', T'')\big[(T'', T')(T', T)\big] = \big[(T''', T'')(T'', T')\big](T', T) = (T''', T).
\]
